\documentclass[../../main.tex]{subfiles}

\begin{document}

\chapter[Landasan Teori dan Notasi Formal]{Landasan Teori dan Notasi Formal}
\label{chap:landasan}

\begin{outcomebox}
\textbf{Sub-CPMK pendukung:} \subcpmkref{1.4}, \subcpmkref{3.1}\\[3pt]
Setelah menyelesaikan bab ini, mahasiswa mampu:
\begin{enumerate}[label=\arabic*.,leftmargin=*]
  \item menggunakan notasi himpunan, relasi, fungsi, dan induksi
        matematika secara tepat;
  \item mendefinisikan alfabet, string, bahasa, dan operasi atas bahasa;
  \item menjelaskan tata bahasa sebagai sistem penulisan ulang dan
        menghasilkan derivasi;
  \item membedakan kelas bahasa pada hierarki Chomsky dan menyiapkan
        data untuk tabel pemetaan Tipe 3--0;
  \item menghasilkan derivasi dan pohon sintaks serta memetakan teori
        ke fase-fase front-end kompilator.
\end{enumerate}
\end{outcomebox}

\section*{Materi Pokok}
\begin{itemize}[leftmargin=*]
  \item Himpunan, relasi, fungsi, dan induksi matematika.
  \item Alfabet, string, bahasa, serta operasi gabungan, konkatenasi,
        dan bintang Kleene.
  \item Tata bahasa sebagai sistem penulisan ulang, derivasi, dan pohon
        urai.
  \item Mesin abstrak: otomata hingga, otomata pushdown, dan mesin
        Turing.
  \item Peta konsep dari teori bahasa ke analisis leksikal, sintaksis,
        dan semantik.
\end{itemize}

\subfile{section-01}
\subfile{section-02}
\subfile{section-03}
\subfile{section-04}
\subfile{section-05}

\section*{Contoh}
\addcontentsline{toc}{section}{Contoh}

\begin{contoh}[Pembuktian dengan induksi]
Buktikan bahwa untuk setiap \(n \ge 1\),
\[
\sum_{i=1}^{n} i = \frac{n(n+1)}{2}.
\]
\textbf{Basis.} Untuk \(n=1\), ruas kiri bernilai \(1\) dan ruas
kanan bernilai \(1(2)/2 = 1\). Basis benar.

\textbf{Langkah induksi.} Asumsikan sifat benar untuk \(k\):
\[
\sum_{i=1}^{k} i = \frac{k(k+1)}{2}.
\]
Maka
\[
\sum_{i=1}^{k+1} i
= \frac{k(k+1)}{2} + (k+1)
= \frac{k(k+1) + 2(k+1)}{2}
= \frac{(k+1)(k+2)}{2}.
\]
Bentuk terakhir adalah rumus untuk \(n = k+1\). Dengan induksi,
sifat terbukti untuk semua \(n \ge 1\).
\end{contoh}

\begin{contoh}[Membuktikan bahasa]
Misalkan \(G = (\{S\}, \{a,b\}, P, S)\) dengan aturan
\(S \to aSb \mid \varepsilon\). Tunjukkan bahwa
\(L(G) = \{a^n b^n \mid n \ge 0\}\).

\textbf{Arah pertama.} Setiap derivasi yang mungkin adalah
\[
S \Rightarrow aSb \Rightarrow a^2 S b^2 \Rightarrow \cdots
\Rightarrow a^n S b^n \Rightarrow a^n b^n.
\]
Dengan induksi pada jumlah langkah, setiap string yang diturunkan
berbentuk \(a^n b^n\).

\textbf{Arah kedua.} Untuk setiap \(n \ge 0\), string \(a^n b^n\)
dapat diturunkan dengan menerapkan aturan \(S \to aSb\) sebanyak
\(n\) kali, lalu \(S \to \varepsilon\). Jadi setiap string berbentuk
tersebut berada di \(L(G)\).

Kedua arah membuktikan \(L(G) = \{a^n b^n \mid n \ge 0\}\).
\end{contoh}

\begin{contoh}[Derivasi dan pohon urai]
Diberikan tata bahasa
\[
E \to E + T \mid T, \quad
T \to T * F \mid F, \quad
F \to (E) \mid \code{id}.
\]
Tentukan derivasi leftmost dan pohon urai untuk
\(\code{id} * (\code{id} + \code{id})\).

Derivasi leftmost:
\[
\begin{aligned}
E &\Rightarrow T\\
  &\Rightarrow T * F\\
  &\Rightarrow F * F\\
  &\Rightarrow \code{id} * F\\
  &\Rightarrow \code{id} * (E)\\
  &\Rightarrow \code{id} * (E + T)\\
  &\Rightarrow \code{id} * (T + T)\\
  &\Rightarrow \code{id} * (F + T)\\
  &\Rightarrow \code{id} * (\code{id} + T)\\
  &\Rightarrow \code{id} * (\code{id} + F)\\
  &\Rightarrow \code{id} * (\code{id} + \code{id}).
\end{aligned}
\]
Pohon uraiya memiliki akar \(E\), cabang kiri \(T\) yang menurun ke
\(\code{id}\), operator \(*\), dan cabang kanan \(F\) yang menurun ke
\((E)\). Di dalam kurung, ekspresi \(E + T\) diurai dengan cara yang
sama seperti contoh sebelumnya.
\end{contoh}

\begin{contoh}[Klasifikasi bahasa]
Tentukan kelas bahasa berikut pada hierarki Chomsky.
\begin{enumerate}[label=(\alph*),leftmargin=*]
  \item \(L_1 = \{w \in \{a,b\}^{*} \mid w \text{ memuat } aba\}\).
  \item \(L_2 = \{a^n b^n \mid n \ge 0\}\).
  \item \(L_3 = \{a^n b^n c^n \mid n \ge 0\}\).
\end{enumerate}
\textbf{Jawaban.}
\begin{enumerate}[label=(\alph*),leftmargin=*]
  \item \(L_1\) regular. Bahasa ini dikenali oleh DFA dengan empat
        state yang mengingat apakah prefiks terakhir adalah
        \(\varepsilon\), \(a\), \(ab\), atau \(aba\).
  \item \(L_2\) bebas konteks, tetapi bukan regular. Tata bahasa
        \(S \to aSb \mid \varepsilon\) membangkitkannya. Bahasa ini
        memerlukan penghitungan jumlah \(a\) dan \(b\) yang tidak
        terbatas, sehingga DFA dengan state berhingga tidak cukup.
  \item \(L_3\) sensitif konteks, tetapi bukan bebas konteks. Bahasa
        ini memerlukan pencocokan tiga kelompok simbol, yang tidak
        dapat dilakukan oleh satu stack.
\end{enumerate}
\end{contoh}

\section*{Aktivitas Pembelajaran dan Praktikum}
\addcontentsline{toc}{section}{Aktivitas Pembelajaran dan Praktikum}

\begin{table}[H]
\centering
\caption{Rencana aktivitas Bab II}
\label{tab:ch02-aktivitas}
\begin{tabularx}{\textwidth}{@{}llX@{}}
\toprule
\textbf{Sesi} & \textbf{Bentuk} & \textbf{Kegiatan} \\
\midrule
1 & Ceramah & Himpunan, relasi, fungsi, dan induksi \\
2 & Latihan & Membuktikan sifat string dan bahasa \\
3 & Lokakarya & Menulis tata bahasa dan derivasi untuk bahasa sederhana \\
4 & Diskusi & Mengklasifikasikan bahasa pada hierarki Chomsky \\
\bottomrule
\end{tabularx}
\end{table}

\subsection*{Lokakarya Notasi}
\begin{enumerate}[label=\arabic*.,leftmargin=*]
  \item Diberikan bahasa \(L = \{w \in \{0,1\}^{*} \mid w \text{
        berakhir dengan } 01\}\), tuliskan lima string anggota dan
        lima string bukan anggota.
  \item Tuliskan tata bahasa bebas konteks untuk \(L\).
  \item Buat derivasi leftmost untuk string \(00101\).
  \item Gambarkan pohon urai untuk derivasi tersebut.
  \item Tentukan kelas bahasa \(L\) pada hierarki Chomsky dan
        jelaskan alasannya.
\end{enumerate}

Bab ini tidak menambahkan increment proyek. Milestone M0 dimulai pada
Bab III setelah arsitektur kompilator dibahas. Namun, mahasiswa
disarankan menyimpan catatan notasi dan contoh derivasi sebagai
rujukan saat menulis spesifikasi Subset-C.

\section*{Latihan dan Refleksi}
\addcontentsline{toc}{section}{Latihan dan Refleksi}

\subsection*{Latihan}
\begin{enumerate}[label=\arabic*.,leftmargin=*]
  \item Buktikan dengan induksi bahwa \(|w^{R}| = |w|\) untuk setiap
        string \(w\).
  \item Diberikan \(L_1 = \{a, aa\}\) dan \(L_2 = \{b, ab\}\).
        Tentukan \(L_1L_2\), \(L_1^{*}\), dan \(L_1 \cup L_2\).
  \item Tuliskan tata bahasa bebas konteks untuk bahasa
        \(\{a^i b^j \mid i \ge j \ge 0\}\).
  \item Buat derivasi leftmost dan pohon urai untuk string
        \(\code{id} + \code{id} * \code{id}\) pada tata bahasa
        aritmetika.
  \item Jelaskan mengapa bahasa \(\{a^n b^n \mid n \ge 0\}\) tidak
        dapat dikenali oleh DFA.
  \item Berikan contoh bahasa sensitif konteks yang bukan bebas
        konteks, lalu jelaskan mengapa satu stack tidak cukup.
  \item Mengapa analisis leksikal cukup memakai otomata hingga,
        sedangkan analisis sintaksis memerlukan otomata pushdown?
  \item Jelaskan hubungan antara pohon urai dan AST. Mengapa AST lebih
        ringkas?
\end{enumerate}

\subsection*{Refleksi}
\begin{reflectionbox}
\begin{enumerate}[label=\arabic*.,leftmargin=*]
  \item Konsep formal mana yang paling sulit Anda pahami? Mengapa?
  \item Apakah Anda dapat menjelaskan perbedaan bahasa regular dan
        bebas konteks tanpa membaca catatan?
  \item Bagaimana notasi formal membantu Anda menulis spesifikasi
        bahasa yang lebih tepat?
\end{enumerate}
\end{reflectionbox}

\section*{Asesmen}
\addcontentsline{toc}{section}{Asesmen}

Asesmen Bab II bersifat formatif dan menjadi bagian dari Kuis Konsep
(3\%). Fokus asesmen adalah ketepatan notasi dan kemampuan
menghubungkan teori dengan fase kompilasi.

\begin{table}[H]
\centering
\caption{Instrumen asesmen Bab II}
\label{tab:ch02-asesmen}
\begin{tabularx}{\textwidth}{@{}lXl@{}}
\toprule
\textbf{Instrumen} & \textbf{Indikator} & \textbf{Bobot} \\
\midrule
Kuis notasi & Ketepatan penggunaan himpunan, string, dan derivasi &
formatif (bagian Kuis Konsep) \\
Tugas derivasi & Derivasi dan pohon urai benar dan lengkap &
formatif (bagian Kuis Konsep) \\
Tugas klasifikasi & Klasifikasi bahasa dan alasan tepat &
formatif (bagian Kuis Konsep) \\
Partisipasi & Keaktifan diskusi dan lokakarya & formatif \\
\bottomrule
\end{tabularx}
\end{table}

Kriteria kelulusan: mahasiswa dapat menulis definisi formal tanpa
kesalahan notasi, menghasilkan derivasi yang benar, dan menjelaskan
hubungan antara kelas bahasa dan mesin pengenal.

\section*{Checklist Pencapaian Kompetensi}
\addcontentsline{toc}{section}{Checklist Pencapaian Kompetensi}

\begin{checklist}
  \item Saya dapat menggunakan notasi himpunan, relasi, dan fungsi.
  \item Saya dapat membuktikan sifat sederhana dengan induksi.
  \item Saya dapat mendefinisikan alfabet, string, dan bahasa.
  \item Saya dapat melakukan operasi gabungan, konkatenasi, dan
        bintang Kleene.
  \item Saya dapat menulis tata bahasa dan menghasilkan derivasi.
  \item Saya dapat menggambar pohon urai dari suatu derivasi.
  \item Saya dapat membedakan bahasa regular, bebas konteks, sensitif
        konteks, dan rekursif terbilang.
  \item Saya dapat menyiapkan data untuk tabel pemetaan Tipe 3--0.
  \item Saya dapat menyiapkan derivasi dan pohon sintaks sebagai
        pendukung Sub-CPMK 3.1.
  \item Saya dapat menjelaskan hubungan teori bahasa dengan fase
        front-end kompilator.
\end{checklist}

\section*{Rangkuman}
\addcontentsline{toc}{section}{Rangkuman}

Pada bab ini telah dibahas:
\begin{itemize}[leftmargin=*]
  \item notasi himpunan, relasi, fungsi, dan induksi matematika;
  \item alfabet, string, bahasa, serta operasi konkatenasi dan
        bintang Kleene;
  \item tata bahasa sebagai sistem penulisan ulang, derivasi, dan
        pohon urai;
  \item hierarki Chomsky dan mesin pengenal yang bersesuaian;
  \item data pendukung tabel pemetaan Tipe 3--0 dan derivasi/pohon;
  \item peta konsep dari teori bahasa ke analisis leksikal,
        sintaksis, dan semantik.
\end{itemize}

Bab III akan memakai seluruh notasi ini untuk menjelaskan arsitektur
kompilator. Kita akan melihat bagaimana string sumber berubah menjadi
token, pohon, dan representasi semantik, serta di mana setiap teknik
yang dipelajari pada bab-bab berikutnya digunakan.

\section*{Evaluasi}
\addcontentsline{toc}{section}{Evaluasi}

\begin{enumerate}[label=\arabic*.,leftmargin=*]
  \item Definisikan tata bahasa \(G = (V,T,P,S)\) dan jelaskan arti
        setiap komponen.
  \item Diberikan \(L = \{w \in \{a,b\}^{*} \mid |w| \text{ genap}\}\).
        Tuliskan tata bahasa regular untuk \(L\).
  \item Jelaskan mengapa \(L = \{a^n b^n \mid n \ge 0\}\) tidak dapat
        dikenali DFA, lalu klasifikasikan bahasa tersebut pada
        hierarki Chomsky.
  \item Buat derivasi leftmost, derivasi rightmost, dan pohon urai
        untuk string \(a + b * c\) pada tata bahasa aritmetika.
  \item Jelaskan perbedaan daya ekspresif otomata hingga, otomata
        pushdown, dan mesin Turing.
  \item Mengapa analisis semantik tidak dapat dilakukan hanya dengan
        tata bahasa bebas konteks? Berikan satu contoh.
\end{enumerate}

\subsection*{Kunci dan Rubrik Singkat}
\begin{itemize}[leftmargin=*]
  \item Soal 1--2: ketepatan definisi dan notasi (70\%), kelengkapan
        (30\%).
  \item Soal 3: ketepatan penjelasan (60\%), klasifikasi (40\%).
  \item Soal 4: kebenaran derivasi (50\%), kebenaran pohon (50\%).
  \item Soal 5--6: ketepatan konsep (60\%), contoh (40\%).
\end{itemize}


\ifSubfilesClassLoaded{%
  \bibliographystyle{apalike}%
  \bibliography{../../references}%
}{}

\end{document}
